\documentclass[../../../include/open-logic-section]{subfiles}

\begin{document}

\olfileid{sth}{z}{milestone}
\olsection{$\Z^-$: Tonggak Wigati}

Kita bakal ngrembug maneh \stagesinf{} ing bagean sabanjure.
Nanging, kanthi Aksioma Tanpa Wates, kita wis tekan
tonggak kang wigati. Saiki kita wis nduweni kabeh aksioma
kang dibutuhake kanggo teori $\Zminus$. Rincine mangkene:

\begin{defn}
Teori $\Zminus$ nduweni aksioma-aksioma iki:
Ekstensionalitas, Gabungan, Pasangan, Himpunan Pangkat,
Tanpa Wates, lan kabeh conto saka skema Pamisahan.
\end{defn}

Jeneng iki nuduhake teori himpunan \emph{Zermelo}
(\emph{dikurangi} sawijining bab kang bakal kita rembug
mengko). Zermelo pantes diajeni amarga sejatine dheweke
ngrumusake teori iki ing tulisane saka taun
\citeyear{Zermelo1908Untersuchungen}.\footnote{Kanggo
cathetan kang narik kawigaten ngenani sejarah lan
rerincen teknise, delengen \citet[Appendix
A]{Potter2004}.}

Teori iki cukup kuwat kanggo nindakake akeh banget
matematika. Mligine, panjenengan \emph{kudu} nyawang maneh
\olref[sfr][][]{part} lan ngyakinake awak dhewe yen kabeh
kang sadurunge kita tindakake kanthi cara naif bisa
ditindakake kanthi luwih formal ing~$\Zminus$. (Sawise
nyoba sawetara wektu, panjenengan bisa langsung maca
\olref[sth][z][arbintersections]{sec}.) Mula wiwit saiki,
tanpa katrangan tambahan, kita bakal nganggep yen
kita nyambut gawe ing $\Zminus$ (saora-orane).

\end{document}
